% from ega4-1-fr.tex
\oldpage[IV-1]{244}

\begin{proposition}[1.9.5]
\label{IV.1.9.5-fr}
Soit $X$ un préschéma.
\begin{enumerate}
  \item[(i)] Pour qu'une partie $E$ de $X$ soit pro-constructible, il faut et
  il suffit que $X-E$ soit ind-constructible.
  \item[(ii)] Toute réunion finie et toute intersection de parties
  pro-constructibles de $X$ sont pro-constructibles. Toute intersection
  finie et toute réunion de parties ind-constructibles de $X$ sont
  ind-constructibles.
  \item[(iii)] Toute intersection (resp. réunion) de parties localement
  constructibles de $X$ est pro-constructible (resp. ind-constructible).
  Inversement, si $X$ est quasi-compact et quasi-séparé, toute partie
  pro-constructible (resp. ind-constructible) de $X$ est intersection (resp.
  réunion) de parties constructibles de $X$.
  \item[(iv)] Soit $(U_\alpha)$ un recouvrement ouvert de $X$. Pour qu'une
  partie $E$ de $X$ soit pro-constructible (resp. ind-constructible) dans
  $X$, il faut et il suffit que pour tout $\alpha$, $E\cap U_\alpha$ soit
  pro-constructible (rep. ind-constructible) dans $U_\alpha$.
  \item[(v)] Toute partie pro-constructible $E$ de $X$ est rétrocompacte
  dans $X$. Pour qu'une partie localement fermée $E$ de $X$ soit
  rétrocompacte dans $X$, il faut et il suffit qu'elle soit
  pro-constructible dans $X$.
  \item[(vi)] Soit $f:X'\to X$ un morphisme de préschémas~; pour toute partie
  pro-constructible (resp. ind-constructible) $E$ de $X$, $f^{-1}(E)$ est
  pro-constructible (resp. ind-constructible) dans $X'$.
  \item[(vii)] Soit $f:X'\to X$ un morphisme quasi-compact~; pour toute
  partie pro-constructible $E'$ de $X'$, $f(E')$ est pro-constructible dans
  $X$~; en particulier $f(X')$ est pro-constructible dans $X$.
  \item[(viii)] Soit $f:X'\to X$ un morphisme localement de présentation
  finie~; pour toute partie ind-constructible $E'$ de $X'$, $f(E')$ est
  ind-constructible dans $X$~; en particulier $f(X')$ est ind-constructible
  dans $X$.
  \item[(ix)] Supposons que $X$ soit quasi-compact et quasi-séparé~; alors,
  pour qu'une partie $E$ de $X$ soit pro-constructible, il faut et il suffit
  qu'il existe un schéma affine $X'$ et un morphisme (nécessairement
  quasi-compact) $f:X'\to X$ tels que $E=f(X')$.
\end{enumerate}
\end{proposition}

\begin{proof}
Les propriétés (i), (ii), (iv) et la première assertion de (iii) résultent
de la définition (1.9.4) et sont valables pour un espace topologique
quelconque, en utilisant la distributivité mutuelle de l'intersection et de
la réunion et le fait que si $T$ est localement constructible dans $X$,
$T\cap U$ est localement constructible dans $U$ pour tout ouvert $U$ de $X$.
Si $X$ est quasi-compact et quasi-séparé et $E$ est pro-constructible dans
$X$, il y a un recouvrement fini $(U_i)$ de $X$ formé d'ouverts affines tels
que $E\cap U_i$ soit intersection de parties constructibles $M_{i\lambda}$
de $U_i$ ($\lambda\in L_i$)~; les $M_{i\lambda}$ sont aussi constructibles
dans $X$ en vertu de (1.2.7) et (0\textsubscript{III}, 9.1.8, (ii)) et $E$ est
intersection des réunions finies $\bigcup_iM_{i,\lambda(i)}$ (où
$\lambda(i)\in L_i$ pour tout $i$), qui sont des parties constructibles dans
$X$~; cela démontre la seconde assertion de (iii). L'assertion (ix) résulte
de (iii) et de (1.9.3). Pour démontrer la première assertion de (v), on peut
se limiter au cas où $X$ est affine, et prouver alors que $E$ est
quasi-compact (0\textsubscript{III}, 9.1.1)~; mais comme $X$ est alors
quasi-séparé, il existe en vertu de (ix) un morphisme quasi-compact
$f:X'\to X$ tel que $E=f(X')$~; comme $X'$ est quasi-compact, il en est de
même de son image $E$ par une application continue.

Pour prouver (vi), on peut se borner, en vertu de (iv), au cas où $X$ est
affine~; la conclusion résulte alors de (iii) et de (1.8.2).

